% ============ 第1章 §1.11 加权范数不等式(原书 pp.88-108) ============
\section{加权范数不等式}\label{sec:1.11}
\index{权}
\index{Ap 类}
\index{加权 Lp 估计!Hardy--Littlewood 极大函数的}
\index{加权 Lp 估计!Hilbert 变换的}
\index{逆 Hölder 不等式}
\index{加权中心极大函数}


在微分方程的边值问题等中,$|x|^{a}$ 形式的权重有时起着重要作用.对这种形式的权重 $w(x)$,Hilbert 变换有时是 $L^{p}(w(x)\dd x)$ 有界的;更一般地,作为使 Hilbert 变换成为 $L^{p}(w(x)\dd x)$ 有界的权重族(类),有 B. Muckenhoupt 发现的 $A_{p}$ 类.本节介绍这个 $A_{p}$ 类.

到前一节为止 Hardy--Littlewood 极大函数也起着重要作用,本节就在与此的关联中展开讨论.本节中也将把函数 $f$ 的 Hardy--Littlewood 极大函数 $M_{1}f$ 简记为 $Mf$.

首先,$\mathbb{R}$ 上的权重是指 $\mathbb{R}$ 上恒不为零,局部可积且非负值的 Lebesgue 可测函数.从下面的定理开始.

\begin{theorem}\label{thm:1.11.1}
设 $1 \leq p < \infty$.设 $\mathbb{R}$ 上的正则正测度 $\mu$ 存在正数 $C_{0}$ 使得
\begin{equation}\tag{1.11.1}
\mu(\{x \in \mathbb{R};\ Mf(x) > \lambda\})^{1/p} \leq \frac{C_{0}}{\lambda}\left(\int_{\mathbb{R}}|f(x)|^{p}\dd\mu(x)\right)^{1/p} \quad (\lambda > 0)
\label{eq:1.11.1}\end{equation}
对一切 $f \in L^{p}(\dd\mu)$ 成立(Hardy--Littlewood 极大函数的弱 $L^{p}(\dd\mu)$ 性).则

(i) $\mu$ 关于 Lebesgue 测度绝对连续(即 $|E| = 0$ 则 $\mu(E) = 0$),存在权重 $w(x)$ 使 $\dd\mu(x) = w(x)\dd x$.且 $w(x) > 0$(a.e. $x \in \mathbb{R}$).

(ii) 对一切开区间 $Q$ 与一切 $f \in L^{p}(\dd\mu)$,
\begin{equation}\tag{1.11.2}
\frac{1}{|Q|}\int_{Q}|f(x)|\dd x \leq C_{0}\left(\frac{1}{\mu(Q)}\int_{Q}|f(x)|^{p}\dd\mu(x)\right)^{1/p}.
\label{eq:1.11.2}\end{equation}

(iii) 对 (i) 中确定的权重 $w$,当 $1 < p < \infty$ 时,
\begin{equation}\tag{1.11.3}
\left(\frac{1}{|Q|}\int_{Q}w(x)\dd x\right)\left(\frac{1}{|Q|}\int_{Q}w(x)^{-1/(p-1)}\dd x\right)^{p-1} \leq C_{0}^{p},
\label{eq:1.11.3}\end{equation}
当 $p = 1$ 时,
\begin{equation}\tag{1.11.4}
\frac{1}{|Q|}\int_{Q}w(x)\dd x \leq C_{0}\operatorname{ess}\inf_{x \in Q}w(x)
\label{eq:1.11.4}\end{equation}
对任意开区间 $Q$ 成立.

(iv) 对每个开区间 $Q$ 与 $Q$ 内的一切 Lebesgue 可测集 $E$,
\begin{equation}\tag{1.11.5}
\frac{|E|}{|Q|} \leq C_{0}\left(\frac{\mu(E)}{\mu(Q)}\right)^{1/p}.
\label{eq:1.11.5}\end{equation}
\end{theorem}

\begin{proof}
(i) 设 $|E| = 0$.为证明 $\mu(E) = 0$,因 $\mu$ 是正则测度,不妨设 $E$ 为紧集.对正数 $\varepsilon$,由 $\mu$ 的正则性,存在开集 $O$ 使 $E \subset O$ 且 $\mu(O \setminus E) = \mu(O) - \mu(E) < \varepsilon$.令 $f(x) = \chi_{O\setminus E}(x)$,则 $f \in L^{p}(\dd\mu)$ 且 $\| f \|_{L^{p}(\dd\mu)} = \left(\int_{\mathbb{R}}\left(\chi_{O\setminus E}(y)\right)^{p}\dd\mu(y)\right)^{1/p} = \mu(O \setminus E)^{1/p} < \varepsilon^{1/p}$.其次,若 $x \in E$,则 $x \in E \subset O$,因 $O$ 是开集,存在开区间 $Q$ 使 $x \in Q \subset O$.因此,因 $|E| = 0$,$|Q|^{-1}\int_{Q}|f(y)|\dd y = |(O \setminus E) \cap Q|/|Q| = |Q \setminus (E \cap Q)|/|Q| = |Q|/|Q| = 1$,故 $Mf(x) = 1$.因此由 \eqref{eq:1.11.1},
\begin{equation*}
\mu(E) \leq \mu(\{x \in \mathbb{R};\ Mf(x) > 1/2\}) \leq 2^{p}C_{0}^{p}\| f \|_{L^{p}(\dd\mu)}^{p} < 2^{p}C_{0}^{p}\varepsilon.
\end{equation*}
因正数 $\varepsilon$ 任意,令 $\varepsilon \to 0$ 得 $\mu(E) = 0$.由此可知 $\mu$ 关于 Lebesgue 测度绝对连续,故存在非负值可测函数 $w(x) \in L_{\mathrm{loc}}^{1}(\mathbb{R})$ 使 $\dd\mu(x) = w(x)\dd x$.为证明 $w(x) > 0$(a.e. $x \in \mathbb{R}$),只需证明对任何具有正测度的可测集 $E$ 有 $\mu(E) > 0$.对使 $E \subset (-b, b)$ 的 $b > 0$,
\begin{equation*}
\frac{1}{2b}\int_{-b}^{b}\chi_{E}(y)\dd y = \frac{|E|}{2b} \leq M\chi_{E}(x) \quad (-b < x < b),
\end{equation*}
故
\begin{equation}\tag{1.11.6}\begin{aligned}
\mu((-b, b)) & \leq \mu(\{x \in \mathbb{R};\ M\chi_{E}(x) > |E|/(4b)\}) \\
& \leq C_{0}^{p}\left(\dfrac{4b}{|E|}\right)^{p}\displaystyle\int_{\mathbb{R}}\left(\chi_{E}(x)\right)^{p}\dd\mu(x) = C_{0}^{p}\left(\frac{4b}{|E|}\right)^{p}\mu(E).
\end{aligned}\label{eq:1.11.6}\end{equation}
因 $\mu$ 是正测度,至少对充分大的 $b$ 有 $\mu((-b, b)) > 0$.因此由 \eqref{eq:1.11.6},$\mu(E) > 0$.

(ii) 设 $Q$ 为开区间,$f \in L^{p}(\dd\mu)$.因 $|Q|^{-1}\int_{Q}|f(y)|\dd y \leq M(f\chi_{Q})(x)$($x \in Q$),
\begin{equation*}
|f|_{Q} := \frac{1}{|Q|}\int_{Q}|f(y)|\dd y \leq \operatorname{essinf}_{x \in Q}M(f\chi_{Q})(x).
\end{equation*}
故由 \eqref{eq:1.11.1},对任意 $0 < a < 1$,
\begin{align*}
\mu(Q) &\leq \mu(\{x \in \mathbb{R};\ M(f\chi_{Q})(x) > a|f|_{Q}\}) \\
\quad &\leq \left(\dfrac{C_{0}}{a|f|_{Q}}\right)^{p}\| f\chi_{Q} \|_{L^{p}(\dd\mu)}^{p} = \left(\dfrac{C_{0}}{a|f|_{Q}}\right)^{p}\displaystyle\int_{Q}|f(y)|^{p}\dd\mu(y).
\end{align*}
由此,因 $0 < a < 1$ 任意,令 $a \to 1$ 得
\begin{equation*}
\frac{1}{|Q|}\int_{Q}|f(x)|\dd x \leq C_{0}\left(\frac{1}{\mu(Q)}\int_{Q}|f(x)|^{p}\dd\mu(x)\right)^{1/p}.
\end{equation*}

(iii) 设 $Q$ 为开区间.首先考虑 $1 < p < \infty$ 的情形.任取正数 $\varepsilon$,则 $(w(x) + \varepsilon)^{-1/(p-1)} \in L^{\infty}(\mathbb{R})$,故
\begin{equation*}
\int_{Q}\left((w(x) + \varepsilon)^{-1/(p-1)}\right)^{p}w(x)\dd x < +\infty.
\end{equation*}
因此在 (ii) 中取 $f = (w(x) + \varepsilon)^{-1/(p-1)}$,
\begin{align*}
\displaystyle\int_{Q}w(x)\dd x\left(\frac{1}{|Q|}\int_{Q}(w(x) + \varepsilon)^{-1/(p-1)}\dd x\right)^{p}&\leq C_{0}^{p}\displaystyle\int_{Q}(w(x) + \varepsilon)^{-p/(p-1)}w(x)\dd x  \\
&= C_{0}^{p}\displaystyle\int_{Q}(w(x) + \varepsilon)^{-1/(p-1)}\frac{w(x)}{w(x) + \varepsilon}\dd x \\
&\leq C_{0}^{p}\displaystyle\int_{Q}(w(x) + \varepsilon)^{-1/(p-1)}\dd x.
\end{align*}
因此
\begin{equation*}
\left(\frac{1}{|Q|}\int_{Q}w(x)\dd x\right)\left(\frac{1}{|Q|}\int_{Q}(w(x) + \varepsilon)^{-1/(p-1)}\dd x\right)^{p-1} \leq C_{0}^{p}.
\end{equation*}
令 $\varepsilon \to 0$,
\begin{equation*}
\left(\frac{1}{|Q|}\int_{Q}w(x)\dd x\right)\left(\frac{1}{|Q|}\int_{Q}w(x)^{-1/(p-1)}\dd x\right)^{p-1} \leq C_{0}^{p}.
\end{equation*}

其次考虑 $p = 1$ 的情形.在 (ii) 中对 $E \subset Q$ 取 $f = \chi_{E}$,
\begin{equation}\tag{1.11.7}
\frac{|E|}{|Q|} \leq C_{0}\frac{\mu(E)}{\mu(Q)} \quad \text{换句话说} \quad \frac{1}{|Q|}\int_{Q}w(x)\dd x \leq C_{0}\frac{1}{|E|}\int_{E}w(x)\dd x.
\label{eq:1.11.7}\end{equation}
这里,若 $a > \operatorname{ess}\inf_{x \in Q}w(x) = \inf\{b > 0;\ |\{x \in Q;\ w(x) < b\}| > 0\}$,则 $|\{x \in Q;\ w(x) < a\}| > 0$,故在 \eqref{eq:1.11.7} 中取 $E = \{x \in Q;\ w(x) < a\}$,
\begin{align*}
\dfrac{1}{|Q|}\displaystyle\int_{Q}w(x)\dd x &\leq C_{0}\frac{1}{|\{y \in Q;\ w(y) < a\}|}\int_{\{y \in Q;\ w(y) < a\}}w(x)\dd x \\
\qquad \qquad \qquad &\leq C_{0}\dfrac{1}{|\{y \in Q;\ w(y) < a\}|}\int_{\{y \in Q;\ w(y) < a\}}a\dd x \leq C_{0}a.
\end{align*}
令 $a \to \operatorname{ess}\inf_{x \in Q}w(x)$,
\begin{equation*}
\frac{1}{|Q|}\int_{Q}w(x)\dd x \leq C_{0}\operatorname{ess}\inf_{x \in Q}w(x).
\end{equation*}

(iv) 对 $E \subset Q$,在 (ii) 中令 $f = \chi_{E}$,
\begin{equation*}
\frac{|E|}{|Q|} \leq C_{0}\left(\frac{\mu(E)}{\mu(Q)}\right)^{1/p}.
\end{equation*}
\end{proof}

后面会证明,由上述 (iii) 可推出 Hardy--Littlewood 极大函数的弱 $L^{p}(w(x)\dd x)$ 性,故引入下面的权重类.

\begin{definition}\label{def:1.11.2}
$\mathbb{R}$ 上的权重 $w(x)$ 满足:对任意区间 $Q$,当 $1 < p < \infty$ 时
\begin{equation}\tag{1.11.8}
\left(\frac{1}{|Q|}\int_{Q}w(x)\dd x\right)\left(\frac{1}{|Q|}\int_{Q}w(x)^{-1/(p-1)}\dd x\right)^{p-1} \leq A^{p},
\label{eq:1.11.8}\end{equation}
当 $p = 1$ 时
\begin{equation}\tag{1.11.9}
\frac{1}{|Q|}\int_{Q}w(x)\dd x \leq A\operatorname{ess}\inf_{x \in Q}w(x)
\label{eq:1.11.9}\end{equation}
若存在满足上述条件的常数 $A$,则称 $w$ 属于(Muckenhoupt 的)$A_{p}$ 类,记为 $w \in A_{p}$.另外,上述常数 $A$ 的下确界称为 $w$ 的 $A_{p}$ 范数,记为 $[w]_{A_{p}}$.

此外,令 $A_{\infty} = \bigcup_{p \geq 1}A_{p}$.
\end{definition}

注意,若权重 $w(x)$ 满足 \eqref{eq:1.11.8} 或 \eqref{eq:1.11.9},则对充分大的区间 $Q$ 有 $\int_{Q}w(x)\dd x > 0$,故由 \eqref{eq:1.11.8} 或 \eqref{eq:1.11.9} 可知 $w(x) > 0$(a.e. $x \in \mathbb{R}$).

另外,把关于 $A_{p}$ 的注意以引理的形式列出.

\begin{lemma}\label{lem:1.11.3}
(i) 若 $1 \leq p < q < \infty$,则 $A_{p} \subset A_{q}$,$[w]_{A_{q}} \leq [w]_{A_{p}}$.

(ii) $[w]_{A_{p}} \geq 1$.

(iii) 若 $1 \leq p < \infty$,$w \in A_{p}$,则对一切开区间 $Q$ 与一切 $f \in L^{p}(\dd\mu)$,
\begin{equation}\tag{1.11.10}
\frac{1}{|Q|}\int_{Q}|f(x)|\dd x \leq [w]_{A_{p}}\left(\frac{1}{w(Q)}\int_{Q}|f(x)|^{p}w(x)\dd x\right)^{1/p}.
\label{eq:1.11.10}\end{equation}

(iv) 设 $1 \leq p < \infty$,$w \in A_{p}$.则对任意开区间 $Q$ 与一切可测集 $E \subset Q$,
\begin{equation*}
\frac{|E|}{|Q|} \leq [w]_{A_{p}}\left(\frac{w(E)}{w(Q)}\right)^{1/p}.
\end{equation*}

(v) 设 $1 \leq p < \infty$,$w \in A_{p}$.又设 $0 < \alpha < 1$.则对任意开区间 $Q$ 与一切满足 $|E| < \alpha|Q|$ 的可测集 $E \subset Q$,存在常数 $0 < \beta < 1$ 使
\begin{equation*}
w(E) \leq \beta w(Q).
\end{equation*}

(vi) 若 $1 \leq p < \infty$,$w \in A_{p}$,则 $w^{-1/(p-1)} \in A_{p'}$($1/p + 1/p' = 1$),且
\begin{equation*}
[w^{-1/(p-1)}]_{A_{p'}} = [w]_{A_{p}}.
\end{equation*}

这里,$w(E) = \int_{E}w(x)\dd x$.
\end{lemma}

\begin{proof}
(i) 设 $w \in A_{p}$,$Q$ 为开区间.当 $p > 1$ 时,因 $(q - 1)/(p - 1) > 1$,由 Hölder 不等式,
\begin{align*}
\left(\dfrac{1}{|Q|}\displaystyle\int_{Q}w(x)^{-1/(q-1)}\dd x\right)^{q-1}&= \left(\dfrac{1}{|Q|}\displaystyle\int_{Q}\left(w(x)^{-1/(p-1)}\right)^{(p-1)/(q-1)}\dd x\right)^{q-1}  \\
&\leq \left(\dfrac{1}{|Q|}\displaystyle\int_{Q}w(x)^{-1/(p-1)}\dd x\right)^{p-1}.
\end{align*}
由此立即得结论.

当 $p = 1$ 时,因 $w(x) \geq [w]_{A_{1}}^{-1}|Q|^{-1}\int_{Q}w(y)\dd y$(a.e. $x \in Q$),故
\begin{align*}
\left(\dfrac{1}{|Q|}\displaystyle\int_{Q}w(x)\dd x\right)\left(\dfrac{1}{|Q|}\int_{Q}w(x)^{-1/(q-1)}\dd x\right)^{q-1}&\leq \left(\dfrac{1}{|Q|}\displaystyle\int_{Q}w(x)\dd x\right)\left(\dfrac{1}{|Q|}\int_{Q}[w]_{A_{1}}^{1/(q-1)}\left(\dfrac{1}{|Q|}\int_{Q}w(x)\dd x\right)^{-1/(q-1)}\dd x\right)^{q-1}  \\
&= [w]_{A_{1}}.
\end{align*}

(ii) 当 $p > 1$ 时,在 \eqref{eq:1.11.8} 中把 $Q$ 缩向 $x$,则对几乎所有的 $x$,$|Q|^{-1}\int_{Q}w(y)\dd y$ 收敛到 $w(x)$,$|Q|^{-1}\int_{Q}w(x)^{-1/(p-1)}\dd x$ 收敛到 $w(x)^{-1/(p-1)}$,故 $1 \leq [w]_{A_{p}}$.当 $p = 1$ 时,注意 $\operatorname{ess}\inf_{x \in Q}w(x) \leq |Q|^{-1}\int_{Q}w(x)\dd x$,则与 $p > 1$ 的情形相同.

(iii) 设 $Q$ 为开区间,$f \in L^{p}(\dd\mu)$.当 $p > 1$ 时,由 Hölder 不等式,
\begin{align*}
\dfrac{1}{|Q|}\displaystyle\int_{Q}|f(x)|\dd x&= \dfrac{1}{|Q|}\displaystyle\int_{Q}|f(x)|w(x)^{1/p}w(x)^{-1/p}\dd x  \\
&\leq \left(\dfrac{1}{|Q|}\displaystyle\int_{Q}|f(x)|^{p}w(x)\dd x\right)^{1/p}\left(\dfrac{1}{|Q|}\int_{Q}w(x)^{-p'/p}\dd x\right)^{1/p'} \\
&= \left(\dfrac{1}{|Q|}\displaystyle\int_{Q}|f(x)|^{p}w(x)\dd x\right)^{1/p}\left(\dfrac{1}{|Q|}\int_{Q}w(x)^{-1/(p-1)}\dd x\right)^{1-1/p} \\
&\leq \left(\dfrac{1}{|Q|}\displaystyle\int_{Q}|f(x)|^{p}w(x)\dd x\right)^{1/p}[w]_{A_{p}}\left(\dfrac{1}{|Q|}\int_{Q}w(x)\dd x\right)^{-1/p} \\
&= [w]_{A_{p}}\left(\dfrac{1}{w(Q)}\displaystyle\int_{Q}|f(x)|^{p}w(x)\dd x\right)^{1/p}.
\end{align*}

当 $p = 1$ 时,利用 $w(x) \geq [w]_{A_{1}}^{-1}|Q|^{-1}\int_{Q}w(y)\dd y$(a.e. $x \in Q$),
\begin{align*}
\dfrac{1}{|Q|}\displaystyle\int_{Q}|f(x)|\dd x & = \dfrac{1}{|Q|}\int_{Q}|f(x)|w(x)w(x)^{-1}\dd x \\
& \leq \dfrac{1}{|Q|}\displaystyle\int_{Q}|f(x)|w(x)\dd x[w]_{A_{1}}\left(\dfrac{1}{|Q|}\int_{Q}w(y)\dd y\right)^{-1} \\
& = [w]_{A_{1}}\dfrac{1}{w(Q)}\displaystyle\int_{Q}|f(x)|w(x)\dd x.
\end{align*}

(iv) 在 (iii) 中取 $f(x) = \chi_{E}$ 即可.

(v) 设 $0 < \alpha < 1$,$Q$ 为开区间,$E$ 为满足 $E \subset Q$ 与 $|E| < \alpha|Q|$ 的可测集.在式~\eqref{eq:1.11.10} 中取 $f = \chi_{Q\setminus E}$,
\begin{align*}
& \left(\dfrac{|Q \setminus E|}{|Q|}\right)^{p} \leq [w]_{A_{p}}^{p}\dfrac{w(Q \setminus E)}{w(Q)} \quad \text{改写为} \\
& \left(1 - \dfrac{|E|}{|Q|}\right)^{p} \leq [w]_{A_{p}}^{p}\left(1 - \dfrac{w(E)}{w(Q)}\right).
\end{align*}
因此由 $|E| < \alpha|Q|$,
\begin{equation*}
(1 - \alpha)^{p} \leq [w]_{A_{p}}^{p}\left(1 - \frac{w(E)}{w(Q)}\right) \quad \text{整理得} \quad \frac{w(E)}{w(Q)} \leq 1 - \frac{(1 - \alpha)^{p}}{[w]_{A_{p}}^{p}}.
\end{equation*}
由 (ii) 知 $[w]_{A_{p}} \geq 1$,又 $0 < \alpha < 1$,故令
\begin{equation*}
\beta = 1 - \frac{(1 - \alpha)^{p}}{[w]_{A_{p}}^{p}},
\end{equation*}
则 $0 < \beta < 1$,满足所求条件.

(vi) 因 $p' = p/(p - 1)$,故 $(p' - 1)(p - 1) = 1$.由它与
\begin{equation}\tag{1.11.11}
\left(\frac{1}{|Q|}\int_{Q}w(x)\dd x\right)\left(\frac{1}{|Q|}\int_{Q}w(x)^{-1/(p-1)}\dd x\right)^{p-1} \leq [w]_{A_{p}}^{p}
\label{eq:1.11.11}\end{equation}
得
\begin{align*}
\left(\dfrac{1}{|Q|}\displaystyle\int_{Q}w(x)^{-1/(p-1)}\dd x\right)\left(\dfrac{1}{|Q|}\int_{Q}\left(w(x)^{-1/(p-1)}\right)^{-1/(p'-1)}\dd x\right)^{p'-1}&\leq [w]_{A_{p}}^{p/(p-1)} = [w]_{A_{p}}^{p'}.
 \end{align*}
因此 $w^{-1/(p-1)} \in A_{p'}$,且 $[w^{-1/(p-1)}]_{A_{p'}} \leq [w]_{A_{p}}$.把 $p$ 与 $p'$ 的角色互换,得 $[w]_{A_{p}} \leq [w^{-1/(p-1)}]_{A_{p'}} \leq [w]_{A_{p}}$,故 $[w^{-1/(p-1)}]_{A_{p'}} = [w]_{A_{p}}$.
\end{proof}

其次,考虑加权中心极大函数,考察其弱 $(1,1)$ 性.

\begin{lemma}\label{lem:1.11.4}
设 $\mu$ 为 $\mathbb{R}$ 上的正测度,满足 $\mu(Q(x, r)) > 0$($x \in \mathbb{R}$,$r > 0$).对 $g \in L_{\mathrm{loc}}^{1}(\mathbb{R}, \dd\mu)$ 定义
\begin{equation*}
\tilde{M}_{\mu}(g)(x) = \sup_{r > 0}\frac{1}{\mu(Q(x, r))}\int_{Q(x, r)}|g(y)|\dd\mu(y).
\end{equation*}
则对一切 $f \in L^{1}(\mathbb{R}, \dd\mu)$,
\begin{equation*}
\mu(\{x \in \mathbb{R};\ \tilde{M}_{\mu}(f)(x) > \lambda\}) \leq \frac{4}{\lambda}\int_{\mathbb{R}}|f(y)|\dd\mu(y) \quad (\lambda > 0).
\end{equation*}
\end{lemma}

\begin{proof}
与定理~\ref{thm:1.5.2} 中 Hardy--Littlewood 中心极大函数 $\tilde{M}f$ 的弱 $(1,1)$ 性的证明完全相同,故省略.
\end{proof}

由上述两个引理,可以证明当 $w \in A_{p}$ 类时 Hardy--Littlewood 极大函数的弱 $L^{p}(w\dd x)$ 性.

\begin{theorem}\label{thm:1.11.5}
设 $1 \leq p < \infty$ 且 $w \in A_{p}$.则对一切 $f \in L^{p}(w(x)\dd x)$,
\begin{equation*}
w(\{x \in \mathbb{R};\ Mf(x) > \lambda\}) \leq \frac{4\cdot 2^{p}[w]_{A_{p}}^{p}}{\lambda^{p}}\int_{\mathbb{R}}|f(x)|^{p}w(x)\dd x \quad (\lambda > 0).
\end{equation*}
\end{theorem}

\begin{proof}
对 $g \in L_{\mathrm{loc}}^{1}(\mathbb{R}, w(x)\dd x)$ 定义
\begin{equation*}
\tilde{M}_{w}g(x) = \sup_{r > 0}\frac{1}{w(Q(x, r))}\int_{Q(x, r)}|g(y)|w(y)\dd y,
\end{equation*}
则由引理~\ref{lem:1.4.4} 与引理~\ref{lem:1.11.3} (iii),
\begin{equation*}
Mf(x) \leq 2[w]_{A_{p}}\left(\tilde{M}_{w}(|f|^{p})(x)\right)^{1/p} \quad (x \in \mathbb{R}).
\end{equation*}
故由引理~\ref{lem:1.11.4},
\begin{align*}
w(\{x \in \mathbb{R};\ Mf(x) > \lambda\}) &\leq w\left(\left\{x \in \mathbb{R};\ \tilde{M}_{w}(|f|^{p})(x) > \lambda^{p}/(2[w]_{A_{p}})^{p}\right\}\right) \\
&\leq \dfrac{4\cdot 2^{p}[w]_{A_{p}}^{p}}{\lambda^{p}}\displaystyle\int_{\mathbb{R}}|f(y)|^{p}w(y)\dd y \quad (\lambda > 0).
\end{align*}
\end{proof}

由这些结果,可以证明对 $A_{p}$ 权重称为逆 Hölder 不等式的下述事实.

\begin{theorem}\label{thm:1.11.6}
设 $1 \leq p < \infty$ 且 $w \in A_{p}$.则存在常数 $C_{0} > 0$,$\gamma > 0$ 使
\begin{equation}\tag{1.11.12}
\left(\frac{1}{|Q|}\int_{Q}w(y)^{1+\gamma}\dd y\right)^{1/(1+\gamma)} \leq \frac{C_{0}}{|Q|}\int_{Q}w(y)\dd y
\label{eq:1.11.12}\end{equation}
对任意区间 $Q$ 成立.这里 $C_{0}, \gamma$ 仅依赖于 $[w]_{A_{p}}$ 与 $p$.
\end{theorem}

\begin{proof}
设 $1 \leq p < \infty$,$w \in A_{p}$,令 $Q = [a, b)$.又设 $0 < \alpha < 1$.进一步令
\begin{equation*}
\alpha_{0} = \frac{1}{|Q|}\int_{Q}w(x)\dd x, \quad \alpha_{k} = \left(\frac{2}{\alpha}\right)^{k}\alpha_{0} \quad (k \in \mathbb{N}_{0}).
\end{equation*}
对 $k \in \mathbb{N}_{0}$,把 $Q$ 作二进分解,对 $w(x)$ 与 $\alpha_{k}$ 作与 Calderón--Zygmund 分解时相同的操作,把所得 $Q$ 的二进分解区间列记为 $\{Q_{k,j}\}$.则

(1) $\alpha_{k} < \dfrac{1}{|Q_{k,j}|}\displaystyle\int_{Q_{k,j}}w(x)\dd x \leq 2\alpha_{k}$.

(2) $w(x) \leq \alpha_{k}$(a.e. $x \in \left(\bigcup_{j}Q_{k,j}\right)^{c}$).

(3) 每个 $Q_{k+1,j}$ 包含在某个 $Q_{k,\ell}$ 中.

(1),(2) 已在 Calderón--Zygmund 分解处证明.关于 (3),因
\begin{equation*}
\frac{1}{|Q_{k+1,j}|}\int_{Q_{k+1,j}}w(x)\dd x > \alpha_{k+1} > \alpha_{k},
\end{equation*}
故在包含 $Q_{k+1,j}$ 的 $Q$ 的二进分解区间 $R$ 中,若不存在满足 $|R|^{-1}\int_{R}w(x)\dd x > \alpha_{k}$ 者,则 $Q_{k+1,j}$ 自身就是某个 $Q_{k,\ell}$;若存在满足 $|R|^{-1}\int_{R}w(x)\dd x > \alpha_{k}$ 者,则其中宽度最大者就是某个 $Q_{k,\ell}$,且 $Q_{k+1,j}$ 包含在其中.

现在,令 $U_{m} = \bigcup_{j}Q_{m,j}$($m \in \mathbb{N}_{0}$),则对 $k \in \mathbb{N}_{0}$,
\begin{align*}
2\alpha_{k} &\geq \dfrac{1}{|Q_{k,\ell}|}\displaystyle\int_{Q_{k,\ell} \cap U_{k+1}}w(y)\dd y \\
&= \dfrac{1}{|Q_{k,\ell}|}\displaystyle\sum_{Q_{k+1,j} \subset Q_{k,\ell}}|Q_{k+1,j}|\frac{1}{|Q_{k+1,j}|}\int_{Q_{k+1,j}}w(y)\dd y \\
&> \dfrac{|Q_{k,\ell} \cap U_{k+1}|}{|Q_{k,\ell}|}\alpha_{k+1} = \dfrac{|Q_{k,\ell} \cap U_{k+1}|}{|Q_{k,\ell}|}\frac{2}{\alpha}\alpha_{k}.
\end{align*}
因此 $|Q_{k,\ell} \cap U_{k+1}| < \alpha|Q_{k,\ell}|$.由引理~\ref{lem:1.11.3} (v),
\begin{equation*}
\frac{w(Q_{k,\ell} \cap U_{k+1})}{w(Q_{k,\ell})} < \beta = 1 - \frac{(1 - \alpha)^{p}}{[w]_{A_{p}}^{p}}.
\end{equation*}
由此,考虑 (3) 得
\begin{align*}
w(U_{k+1}) &= w\left(\left(\bigcup_{\ell}Q_{k,\ell}\right) \cap U_{k+1}\right) \\
&= \displaystyle\sum_{\ell}w(Q_{k,\ell} \cap U_{k+1}) \\
&< \displaystyle\sum_{\ell}\beta w(Q_{k,\ell}) = \beta w(U_{k}).
\end{align*}
故 $w(U_{k}) < \beta^{k}w(U_{0})$.另外,同样可得 $|U_{k}| < \alpha^{k}|U_{0}|$,故 $|U_{k}| \to 0$($k \to \infty$).因此 $|\bigcap_{k}U_{k}| = 0$.由此
\begin{equation}\tag{1.11.13}
\left|Q \setminus \left((Q \setminus U_{0}) \cup \left(\bigcup_{k=0}^{\infty}(U_{k} \setminus U_{k+1})\right)\right)\right| = 0.
\label{eq:1.11.13}\end{equation}
这里,取正数 $\gamma$ 使
\begin{equation*}
(2\alpha^{-1})^{\gamma}\beta < 1.
\end{equation*}
则由 \eqref{eq:1.11.13} 与 (2),
\begin{align*}
\displaystyle\int_{Q}w(y)^{1+\gamma}\dd y &= \int_{Q \setminus U_{0}}w(y)^{\gamma}w(y)\dd y + \sum_{k=0}^{\infty}\int_{U_{k} \setminus U_{k+1}}w(y)^{\gamma}w(y)\dd y \\
\qquad &\leq \alpha_{0}^{\gamma}w(Q \setminus U_{0}) + \displaystyle\sum_{k=0}^{\infty}\alpha_{k+1}^{\gamma}w(U_{k}) \\
\qquad &\leq \alpha_{0}^{\gamma}w(Q \setminus U_{0}) + \displaystyle\sum_{k=0}^{\infty}\left((2\alpha^{-1})^{k+1}\alpha_{0}\right)^{\gamma}\beta^{k}w(U_{0}) \\
\qquad &\leq \alpha_{0}^{\gamma}\left(1 + (2\alpha^{-1})^{\gamma}\sum_{k=0}^{\infty}(2\alpha^{-1})^{\gamma k}\beta^{k}\right)\left(w(Q \setminus U_{0}) + w(U_{0})\right) \\
\qquad &= \left(\dfrac{1}{|Q|}\displaystyle\int_{Q}w(y)\dd y\right)^{\gamma}\left(1 + \dfrac{(2\alpha^{-1})^{\gamma}}{1 - (2\alpha^{-1})^{\gamma}\beta}\right)\int_{Q}w(y)\dd y.
\end{align*}
因此,选取正数 $\gamma, C$ 满足
\begin{align*}
\gamma &< \dfrac{-\log\beta}{\log 2 - \log\alpha} = \dfrac{\log[w]_{A_{p}}^{p} - \log\left([w]_{A_{p}}^{p} - (1 - \alpha)^{p}\right)}{\log 2 - \log\alpha}, \\
C^{1+\gamma} &= 1 + \dfrac{(2\alpha^{-1})^{\gamma}}{1 - (2\alpha^{-1})^{\gamma}\beta} = 1 + \dfrac{(2\alpha^{-1})^{\gamma}}{1 - (2\alpha^{-1})^{\gamma}\left(1 - (1 - \alpha)^{p}/[w]_{A_{p}}^{p}\right)},
\end{align*}
则
\begin{equation*}
\left(\frac{1}{|Q|}\int_{Q}w(y)^{1+\gamma}\dd y\right)^{1/(1+\gamma)} \leq \frac{C}{|Q|}\int_{Q}w(y)\dd y.
\end{equation*}
\end{proof}

作为逆 Hölder 不等式的直接推论,

\begin{lemma}\label{lem:1.11.7}
(i) 设 $1 \leq p < \infty$,$w \in A_{p}$,$\gamma, C_{0}$ 为定理~\ref{thm:1.11.6} 中的常数.则对一切区间 $Q$ 与任意可测集 $E \subset Q$,
\begin{equation*}
\frac{w(E)}{w(Q)} \leq C_{0}\left(\frac{|E|}{|Q|}\right)^{\gamma/(\gamma+1)}.
\end{equation*}

(ii) 设 $1 < p < \infty$,$w \in A_{p}$.则存在仅与 $[w]_{A_{p}}$ 和 $p$ 有关的常数 $p_{0}$($1 < p_{0} < p$)使 $w \in A_{p_{0}}$.
\end{lemma}

\begin{proof}
(i) 设 $Q$ 为区间,$E$ 为其可测子集.由 Hölder 不等式与定理~\ref{thm:1.11.6},
\begin{align*}
\displaystyle\int_{E}w(x)\dd x &\leq \left(\int_{E}\dd x\right)^{\gamma/(\gamma+1)}\left(\int_{E}w(x)^{1+\gamma}\dd x\right)^{1/(1+\gamma)} \\
&\leq |E|^{\gamma/(\gamma+1)}|Q|^{1/(1+\gamma)}\dfrac{C_{0}}{|Q|}\int_{Q}w(x)\dd x.
\end{align*}
整理得
\begin{equation*}
\frac{w(E)}{w(Q)} \leq C_{0}\left(\frac{|E|}{|Q|}\right)^{\gamma/(1+\gamma)}.
\end{equation*}

(ii) 由引理~\ref{lem:1.11.3} (vi),$w^{-1/(p-1)} \in A_{p'}$($1/p + 1/p' = 1$)且 $[w^{-1/(p-1)}]_{A_{p'}} = [w]_{A_{p}}$.对这个 $w^{-1/(p-1)}$ 使用逆 Hölder 不等式(定理~\ref{thm:1.11.6}),把此时的常数记为 $\gamma$,$C_{0}$,用 $(1 + \gamma)/(p - 1) = 1/(p_{0} - 1)$ 确定 $p_{0}$,则 $1 < p_{0} < p$ 且
\begin{equation*}
\begin{aligned}
\left(\dfrac{1}{|Q|}\displaystyle\int_{Q}w(x)^{-1/(p_{0}-1)}\dd x\right)^{p_{0}-1}&= \left(\dfrac{1}{|Q|}\displaystyle\int_{Q}\left(w(x)^{-1/(p-1)}\right)^{1+\gamma}\dd x\right)^{(p-1)/(1+\gamma)}
 \end{aligned}
\end{equation*}
\begin{align*}
&\leq \left(C_{0}\dfrac{1}{|Q|}\displaystyle\int_{Q}w(x)^{-1/(p-1)}\dd x\right)^{p-1} \\
&\leq C_{0}^{p-1}[w]_{A_{p}}^{p}\left(\dfrac{1}{|Q|}\displaystyle\int_{Q}w(x)\dd x\right)^{-1}.
\end{align*}
也就是说,
\begin{equation*}
w \in A_{p_{0}}, \quad [w]_{A_{p_{0}}} \leq C_{0}^{(p-1)/p_{0}}[w]_{A_{p}}^{p/p_{0}}.
\end{equation*}
\end{proof}

作为这个应用,当 $1 < p < \infty$ 时,得到强化定理~\ref{thm:1.11.5} 的 Hardy--Littlewood 极大函数的加权 $L^{p}$ 估计.

\begin{lemma}\label{lem:1.11.8}
设 $1 < p < \infty$,$w \in A_{p}$.则存在常数 $C$ 使
\begin{equation*}
\| Mf \|_{L^{p}(w\dd x)} \leq C\| f \|_{L^{p}(w\dd x)} \quad (f \in L^{p}(w\dd x)).
\end{equation*}
\end{lemma}

\begin{proof}
设 $1 < p < \infty$,$w \in A_{p}$,则由引理~\ref{lem:1.11.7} (ii),存在 $p_{0}$($1 < p_{0} < p$)使 $w \in A_{p_{0}}$.因此由定理~\ref{thm:1.11.5},存在常数 $C_{0}$ 使
\begin{align*}
w(\{x \in \mathbb{R};\ Mf(x) > \lambda\}) & \leq \dfrac{C_{0}}{\lambda^{p_{0}}}\displaystyle\int_{\mathbb{R}}|f(x)|^{p_{0}}w(x)\dd x \\
& \quad (f \in L^{p_{0}}(w\dd x),\ \lambda > 0).
\end{align*}
另外,显然
\begin{equation*}
\| Mf \|_{L^{\infty}(w\dd x)} \leq \| f \|_{L^{\infty}(w\dd x)}.
\end{equation*}
所以在 $p_{0}$ 与 $\infty$ 之间应用 Marcinkiewicz 插值定理~\ref{thm:1.8.6},
\begin{equation*}
\| Mf \|_{L^{p}(w\dd x)} \leq C\| f \|_{L^{p}(w\dd x)} \quad (f \in L^{p}(w\dd x)).
\end{equation*}
\end{proof}

另外,可以证明 Hardy--Littlewood 极大函数与锐极大函数之间的加权``good $\lambda$-不等式'',从而也得到 Hilbert 变换的加权 $L^{p}$ 结论.

\begin{lemma}\label{lem:1.11.9}
设 $1 \leq p < \infty$,$w \in A_{p}$.则当 $f \in L_{\mathrm{loc}}^{1}(\mathbb{R})$ 时,对任意 $\lambda > 0$,$\gamma > 0$,存在常数 $C > 0$ 与 $\delta > 0$ 使
\begin{equation*}
w(\{x \in \mathbb{R};\ Mf(x) > 2\lambda,\ f^{\#}(x) < \gamma\lambda\}) \leq C\gamma^{\delta}w(\{x \in \mathbb{R};\ Mf(x) > \lambda\}).
\end{equation*}
\end{lemma}

\begin{proof}
证明只需对无权重情形的引理~\ref{lem:1.6.5} 的证明作部分修改即可.到式~\eqref{eq:1.6.1},即
\begin{equation*}
|\{x \in Q;\ Mf(x) > 2\lambda,\ f^{\#}(x) < \gamma\lambda\}| \leq 96\gamma|Q|
\end{equation*}
为止完全相同.由这个不等式与引理~\ref{lem:1.11.7} (i)(把那里的常数 $C_{0}$,$\gamma$ 记为 $C_{0}$,$\beta$),得
\begin{equation}\tag{1.11.14}
w(\{x \in Q;\ Mf(x) > 2\lambda,\ f^{\#}(x) < \gamma\lambda\}) \leq C_{0}(96\gamma)^{\beta/(1+\beta)}w(Q).
\label{eq:1.11.14}\end{equation}
因此,把对 $Q_{j}$ 的估计式~\eqref{eq:1.11.14} 对 $Q_{j}$ 求和,因 $Q_{j}$ 的内部互不相交,
\begin{align*}
w(\{x \in \mathbb{R};\ Mf(x) > 2\lambda,\ f^{\#}(x) < \gamma\lambda\})&= w(\{x \in E_{\lambda};\ Mf(x) > 2\lambda,\ f^{\#}(x) < \gamma\lambda\})  \\
&\leq C_{0}(96\gamma)^{\beta/(1+\beta)}w(E_{\lambda}) = C_{0}(96\gamma)^{\beta/(1+\beta)}w(\{x \in \mathbb{R};\ Mf(x) > \lambda\}).
\end{align*}
\end{proof}

把这个加权的``good $\lambda$-不等式''与引理~\ref{lem:1.6.6} 组合,就得到定理~\ref{thm:1.6.8} 的加权版本.

\begin{lemma}\label{lem:1.11.10}
设 $1 \leq p < \infty$,$w \in A_{p}$.则当 $Mf \in L^{p}(w\dd x)$ 时,存在常数 $C$ 使
\begin{equation*}
\| Mf \|_{L^{p}(w\dd x)} \leq C\| f^{\#} \|_{L^{p}(w\dd x)} \quad (f \in L^{p}(w\dd x)).
\end{equation*}
\end{lemma}

\begin{proof}
只需对无权重情形的引理~\ref{lem:1.6.6} 的证明稍作修改即可,故省略证明.
\end{proof}

有了以上准备,就能给出 Hilbert 变换的加权 $L^{p}$ 估计.

\begin{theorem}\label{thm:1.11.11}
(i) 当 $1 < p < \infty$,$w \in A_{p}$ 时,存在常数 $C$ 使
\begin{equation*}
\| Hf \|_{L^{p}(w\dd x)} \leq C\| f \|_{L^{p}(w\dd x)} \quad (f \in L^{p}(w\dd x)).
\end{equation*}

(ii) 当 $w \in A_{1}$ 时,存在常数 $C$ 使
\begin{equation*}
w(\{x \in \mathbb{R};\ |Hf(x)| > \lambda\}) \leq \frac{C}{\lambda}\| f \|_{L^{1}(w\dd x)} \quad (f \in L^{1}(w\dd x)).
\end{equation*}
\end{theorem}

\begin{proof}
(i) 设 $f \in \mathcal{S}(\mathbb{R})$.已知 $|Hf(x)| \leq C_{1}(1 + |x|)^{-1}$,故 $M(Hf)(x) \leq 4C_{1}(1 + |x|)^{-1}\log(e + |x|)$(这个计算留给读者).于是 $M(Hf) \in L^{p}(w\dd x)$.实际上,由引理~\ref{lem:1.11.7} (ii) 存在 $1 < p_{0} < p$ 使 $w \in A_{p_{0}}$,故利用引理~\ref{lem:1.11.3} (iv),对 $j \geq 1$ 可导出 $w(\{|x| < 2^{j+1}\})/w(\{|x| < 2\}) \leq [w]_{A_{p_{0}}}^{p_{0}}(2^{j+1}/2)^{p_{0}} \leq [w]_{A_{p}}^{p_{0}}2^{jp_{0}}$,故当 $j \geq 1$ 时,
\begin{align*}
\displaystyle\int_{2^{j} \leq |x| < 2^{j+1}}M(Hf)(x)^{p}w(x)\dd x&\leq \left(4C_{1}\dfrac{\log(e + 2^{j+1})}{1 + 2^{j}}\right)^{p}\displaystyle\int_{2^{j} \leq |x| < 2^{j+1}}w(x)\dd x  \\
&\leq (4C_{1}\log 2)^{p}(j + 2)^{p}2^{-jp}\displaystyle\int_{|x| < 2^{j+1}}w(x)\dd x \\
&= (4C_{1}\log 2)^{p}(j + 2)^{p}2^{-jp}\dfrac{w(\{|x| < 2^{j+1}\})}{w(\{|x| < 2\})}w(\{|x| < 2\}) \\
&\leq (4C_{1}\log 2)^{p}(j + 2)^{p}2^{-jp}C_{2}2^{jp_{0}}w(\{|x| < 2\}) \\
&\leq C_{3}j^{p}2^{-(p-p_{0})j}w(\{|x| < 2\}).
\end{align*}
因此,
\begin{equation*}
\int_{\mathbb{R}}M(Hf)(x)^{p}w(x)\dd x
\end{equation*}
\begin{align*}
&= \displaystyle\int_{|x| < 2}M(Hf)(x)^{p}w(x)\dd x + \sum_{j=1}^{\infty}\int_{2^{j} \leq |x| < 2^{j+1}}M(Hf)(x)^{p}w(x)\dd x \\
&\leq \left(4C_{1}\log(e + 2)\right)^{p}w(\{|x| < 2\}) + \displaystyle\sum_{j=1}^{\infty}C_{3}j^{p}2^{-(p-p_{0})j}w(\{|x| < 2\}) < +\infty.
\end{align*}
由此利用引理~\ref{lem:1.11.10},
\begin{equation}\tag{1.11.15}
\| Hf \|_{L^{p}(w\dd x)} \leq \| M(Hf) \|_{L^{p}(w\dd x)} \leq C_{4}\| (Hf)^{\#} \|_{L^{p}(w\dd x)}.
\label{eq:1.11.15}\end{equation}
另一方面,由引理~\ref{lem:1.11.7} (ii) 存在 $1 < p_{0} < p$ 使 $w \in A_{p_{0}}$,又由系~\ref{cor:1.7.4},
\begin{equation*}
(Hf)^{\#}(x) \leq C_{5}M_{p/p_{0}}f(x) \quad (x \in \mathbb{R}),
\end{equation*}
故由引理~\ref{lem:1.11.8},
\begin{equation}\tag{1.11.16}\begin{aligned}
\| (Hf)^{\#} \|_{L^{p}(w\dd x)} & \leq C_{5}\| M_{p/p_{0}}f \|_{L^{p}(w\dd x)} \\
& = C_{5}\left(\displaystyle\int_{\mathbb{R}}\left(M(|f|^{p/p_{0}})(x)\right)^{p_{0}}w(x)\dd x\right)^{1/p} \\
& \leq C_{5}\left[C_{6}\left(\displaystyle\int_{\mathbb{R}}\left(|f|^{p/p_{0}}(x)\right)^{p_{0}}w(x)\dd x\right)^{1/p_{0}}\right]^{p_{0}/p} \\
& \leq C_{6}\| f \|_{L^{p}(w\dd x)}.
\end{aligned}\label{eq:1.11.16}\end{equation}
合并 \eqref{eq:1.11.15},\eqref{eq:1.11.16},
\begin{equation*}
\| Hf \|_{L^{p}(w\dd x)} \leq C_{4}C_{6}\| f \|_{L^{p}(w\dd x)}.
\end{equation*}
可以证明 $\mathcal{S}(\mathbb{R})$ 在 $L^{p}(w\dd x)$ 中稠密,故得到所求结论.

(ii) 设 $w \in A_{1}$,$f \in L^{1}(w\dd x) \cap L^{2}(w\dd x)$,$\lambda > 0$.首先确认 $f \in L_{\mathrm{loc}}^{1}(\mathbb{R})$.设 $Q$ 为区间,则
\begin{align}
\frac{1}{|Q|}\int_{Q}|f(x)|\dd x &\leq \frac{1}{|Q|}\int_{Q}|f(x)|w(x)\dd x\cdot\frac{1}{\operatorname{essinf}_{x \in Q}w(x)} \notag \\
&\leq \frac{1}{|Q|}\int_{Q}|f(x)|w(x)\dd x\cdot\frac{[w]_{A_{1}}}{\frac{1}{|Q|}\int_{Q}w(x)\dd x} \notag \\
&= \frac{[w]_{A_{1}}}{w(Q)}\int_{Q}|f(x)|w(x)\dd x \tag{1.11.17}\label{eq:1.11.17} \\
&\leq \frac{[w]_{A_{1}}}{w(Q)}\| f \|_{L^{1}(w\dd x)} < +\infty. \tag{1.11.18}\label{eq:1.11.18}
\end{align}
因此 $f \in L_{\mathrm{loc}}^{1}(\mathbb{R})$.其次,对 Calderón--Zygmund 分解的引理~\ref{lem:1.8.1} 的证明稍作修改,确认现在的情形下也能作同样的分解.引入二进极大函数
\begin{equation*}
M_{d}f(x) = \sup_{x \in Q \in \mathbb{D}}\frac{1}{|Q|}\int_{Q}|f(y)|\dd y,
\end{equation*}
令
\begin{equation*}
\Omega = \{x \in \mathbb{R};\ M_{d}f(x) > \lambda\}.
\end{equation*}
若 $x \in \Omega$,则由上确界的性质,存在包含 $x$ 且满足 $1/|Q|\int_{Q}|f(y)|\dd y > \lambda$ 的二进区间 $Q$.反之,若 $1/|Q|\int_{Q}|f(y)|\dd y > \lambda$,则对所有 $x \in Q$ 有 $M_{d}f(x) > \lambda$.因此,把满足 $1/|Q|\int_{Q}|f(y)|\dd y > \lambda$ 的 $Q \in \mathbb{D}$ 的全体记为 $\{R_{j}\}$,则
\begin{equation*}
\Omega = \bigcup_{j}R_{j}.
\end{equation*}
对 $\{R_{j}\}$ 中不断扩大的子列 $S_{1} \subset S_{2} \subset \cdots$,由 \eqref{eq:1.11.18} 与引理~\ref{lem:1.11.7},
\begin{align*}
\dfrac{1}{|S_{j}|}\displaystyle\int_{S_{j}}|f(x)|\dd x & \leq [w]_{A_{1}}\dfrac{w(S_{1})}{w(S_{j})}\dfrac{\| f \|_{L^{1}(w\dd x)}}{w(S_{1})} \\
& \leq C_{0}[w]_{A_{1}}\left(\dfrac{|S_{1}|}{|S_{j}|}\right)^{\gamma/(1+\gamma)}\dfrac{\| f \|_{L^{1}(w\dd x)}}{w(S_{1})},
\end{align*}
故可知这种子列只有有限个.因此由引理~\ref{lem:1.6.3},在 $\{R_{j}\}$ 中取极大二进区间 $\{Q_{j}\}$,则
\begin{equation*}
\Omega = \bigcup_{j}Q_{j}, \quad Q_{k} \cap Q_{\ell} = \emptyset \quad (k \neq \ell),
\end{equation*}
且由 $Q_{j}$ 的极大性,
\begin{equation*}
\lambda < \frac{1}{|Q_{j}|}\int_{Q_{j}}|f(x)|\dd x = |f|_{Q_{j}} \leq 2\lambda.
\end{equation*}
也利用 \eqref{eq:1.11.17},
\begin{equation*}
\lambda < \frac{[w]_{A_{1}}}{w(Q_{j})}\int_{Q_{j}}|f(x)|w(x)\dd x,
\end{equation*}
故
\begin{equation}\tag{1.11.19}
w(\Omega) = \sum_{j}w(Q_{j}) < \sum_{j}\frac{[w]_{A_{1}}}{\lambda}\int_{Q_{j}}|f(x)|w(x)\dd x = \frac{[w]_{A_{1}}}{\lambda}\| f \|_{L^{1}(w\dd x)}.
\label{eq:1.11.19}\end{equation}
另外,由 $M_{d}f(x) \leq \lambda$($x \in \Omega^{c}$),
\begin{equation*}
|f(x)| \leq \lambda \quad (\text{a.e. } x \in \Omega^{c}).
\end{equation*}
于是,像定理~\ref{thm:1.8.2} 的证明那样定义
\begin{align*}
g(x) &= \begin{cases} f(x) & (x \in \Omega^{c}) \\ \dfrac{1}{|Q_{j}|}\displaystyle\int_{Q_{j}}f(x)\dd x & (x \in Q_{j},\ j = 1, 2, \dots) \end{cases} \\[10pt]
f_{j}(x) &= \left(f(x) - \dfrac{1}{|Q_{j}|}\displaystyle\int_{Q_{j}}f(y)\dd y\right)\chi_{Q_{j}} \\[8pt]
b(x) &= \displaystyle\sum_{j}f_{j}(x).
\end{align*}
则 $f(x) = g(x) + b(x)$,且
\begin{equation*}
|g(x)| \leq 2\lambda \quad (\text{a.e. } x \in \mathbb{R}).
\end{equation*}
进而,由 \eqref{eq:1.11.17} 与 $[w]_{A_{1}} \geq 1$,
\begin{align*}
\displaystyle\int_{\mathbb{R}}|g(x)|w(x)\dd x&\leq \displaystyle\int_{\Omega^{c}}|f(x)|w(x)\dd x + \sum_{j}\int_{Q_{j}}\left(\frac{1}{|Q_{j}|}\int_{Q_{j}}|f(y)|\dd y\right)w(x)\dd x  \\
&= \displaystyle\int_{\Omega^{c}}|f(x)|w(x)\dd x + \sum_{j}w(Q_{j})\frac{1}{|Q_{j}|}\int_{Q_{j}}|f(y)|\dd y \\
&\leq \displaystyle\int_{\Omega^{c}}|f(x)|w(x)\dd x + \sum_{j}[w]_{A_{1}}\int_{Q_{j}}|f(y)|w(y)\dd y \\
&\leq [w]_{A_{1}}\displaystyle\int_{\mathbb{R}}|f(x)|w(x)\dd x.
\end{align*}
由此 $g \in L^{1}(w\dd x) \cap L^{\infty}(w\dd x)$,故亦有 $g \in L^{1}(w\dd x) \cap L^{2}(w\dd x)$.

其次,设 $Q_{j}$ 的中心为 $x_{j}$,
\begin{equation*}
Q_{j}^{*} = 2Q_{j} = Q(x_{j}, 2\ell(Q_{j})), \quad \Omega^{*} = \bigcup_{j}Q_{j}^{*},
\end{equation*}
则由引理~\ref{lem:1.11.3} (iv) 与 \eqref{eq:1.11.19},令 $C_{1} = 2[w]_{A_{1}}$,
\begin{equation}\tag{1.11.20}
w(\Omega^{*}) \leq \sum_{j}w(Q_{j}^{*}) = \sum_{j}C_{1}w(Q_{j}) = C_{1}w(\Omega) \leq \frac{C_{2}}{\lambda}\| f \|_{L^{1}(w\dd x)}.
\label{eq:1.11.20}\end{equation}
现在,由 $f(x) = g(x) + b(x)$ 得 $Hf(x) = Hg(x) + Hb(x)$,若 $|Hg(x)| \leq \lambda/2$ 且 $|Hb(x)| \leq \lambda/2$,则 $|Hf(x)| \leq |Hg(x)| + |Hb(x)| \leq \lambda$,故
\begin{align*}
& \{x \in \mathbb{R};\ |Hf(x)| > \lambda,\ x \in (\Omega^{*})^{c}\} \\
& \qquad \subset \{x \in (\Omega^{*})^{c};\ |Hg(x)| > \lambda/2\} \cup \{x \in (\Omega^{*})^{c};\ |Hb(x)| > \lambda/2\}.
\end{align*}
令 $E_{1} = \{x \in (\Omega^{*})^{c};\ |Hg(x)| > \lambda/2\}$,$E_{2} = \{x \in (\Omega^{*})^{c};\ |Hb(x)| > \lambda/2\}$.首先,利用 $H$ 的 $L^{2}(w\dd x)$ 有界性与 \eqref{eq:1.11.19},
\begin{equation}\tag{1.11.21}\begin{aligned}
w(E_{1}) &\leq \dfrac{1}{(\lambda/2)^{2}}\displaystyle\int_{E_{1}}|Hg(x)|^{2}w(x)\dd x \leq \dfrac{4}{\lambda^{2}}\int_{\mathbb{R}}|Hg(x)|^{2}w(x)\dd x \\
&\leq \dfrac{C_{3}}{\lambda^{2}}\displaystyle\int_{\mathbb{R}}|g(x)|^{2}w(x)\dd x \\
&= \dfrac{C_{3}}{\lambda^{2}}\left(\displaystyle\int_{\Omega}|g(x)|^{2}w(x)\dd x + \int_{\Omega^{c}}|g(x)|^{2}w(x)\dd x\right) \\
&= \dfrac{C_{3}}{\lambda^{2}}\left(\displaystyle\sum_{j}\int_{Q_{j}}|f_{Q_{j}}|^{2}w(x)\dd x + \lambda\int_{\Omega^{c}}|g(x)|w(x)\dd x\right) \\
&\quad \quad \quad \quad \quad \quad \quad \quad \quad \quad \quad \quad \quad \quad \quad (|g(x)| \leq \lambda,\ \text{a.e. } x \in \Omega^{c}) \\
&\leq \dfrac{C_{3}}{\lambda^{2}}\left(\displaystyle\sum_{j}w(Q_{j})(2\lambda)^{2} + \lambda\int_{\mathbb{R}}|f(x)|w(x)\dd x\right) \\
&\quad \quad \quad \quad \quad \quad \quad \quad \quad \quad (|f_{Q_{j}}| \leq 2\lambda,\ g(x) = f(x)\ (x \in \Omega^{c})) \\
&= 4C_{3}\displaystyle\sum_{j}w(Q_{j}) + \dfrac{C_{3}}{\lambda}\| f \|_{L^{1}(w\dd x)} \leq \dfrac{C_{4}}{\lambda}\| f \|_{L^{1}(w\dd x)}.
\end{aligned}\label{eq:1.11.21}\end{equation}
其次,若 $x \in (Q_{j}^{*})^{c}$,则
\begin{equation*}
|Hb(x)| \leq \sum_{j}|Hf_{j}(x)| \leq \frac{8\lambda}{\pi}\sum_{j}\frac{\ell(Q_{j})^{2}}{|x - x_{j}|^{2} + \ell(Q_{j})^{2}},
\end{equation*}
故注意 $\ell(Q_{j}) = |Q_{j}|$,
\begin{equation}\tag{1.11.22}\begin{aligned}
w(E_{2}) &\leq \dfrac{1}{\lambda/2}\displaystyle\int_{E_{2}}|Hb(x)|w(x)\dd x \leq \dfrac{2}{\lambda}\int_{(\Omega^{*})^{c}}|Hb(x)|w(x)\dd x \\
\quad &\leq \dfrac{16}{\pi}\displaystyle\int_{(\Omega^{*})^{c}}\sum_{j}\frac{|Q_{j}|^{2}}{|x - x_{j}|^{2} + |Q_{j}|^{2}}w(x)\dd x \\
\quad &\leq \dfrac{16}{\pi}\displaystyle\sum_{j}|Q_{j}|\int_{-\infty}^{\infty}\frac{|Q_{j}|}{|x - x_{j}|^{2} + |Q_{j}|^{2}}w(x)\dd x.
\end{aligned}\label{eq:1.11.22}\end{equation}
这里,令 $2^{\ell}Q_{j} = \left[x_{j} - 2^{\ell-1}|Q_{j}|,\ x_{j} + 2^{\ell-1}|Q_{j}|\right)$,则由 $w \in A_{1}$,
\begin{align*}
& \displaystyle\int_{-\infty}^{\infty}\frac{|Q_{j}|}{|x - x_{j}|^{2} + |Q_{j}|^{2}}w(x)\dd x \\
& = \displaystyle\int_{|x - x_{j}| < |Q_{j}|/2}\frac{|Q_{j}|}{|x - x_{j}|^{2} + |Q_{j}|^{2}}w(x)\dd x \\
& \qquad + \displaystyle\sum_{m=0}^{\infty}\int_{2^{m-1}|Q_{j}| \leq |x - x_{j}| < 2^{m}|Q_{j}|}\frac{|Q_{j}|}{|x - x_{j}|^{2} + |Q_{j}|^{2}}w(x)\dd x \\
& \leq \dfrac{1}{|Q_{j}|}\displaystyle\int_{|x - x_{j}| < |Q_{j}|/2}w(x)\dd x \\
& \qquad + \displaystyle\sum_{m=0}^{\infty}\frac{1}{2^{m-1}}\int_{|x - x_{j}| < 2^{m}|Q_{j}|}\frac{1}{2^{m-1}|Q_{j}|}w(x)\dd x \\
& \leq \dfrac{w(Q_{j})}{|Q_{j}|} + \displaystyle\sum_{m=0}^{\infty}\frac{1}{2^{m-2}}\frac{w(2^{m}Q_{j})}{|2^{m}Q_{j}|} \\
& \leq \dfrac{w(Q_{j})}{|Q_{j}|} + \displaystyle\sum_{m=0}^{\infty}\frac{1}{2^{m-2}}[w]_{A_{1}}\operatorname{essinf}_{x \in 2^{m}Q_{j}}w(x) \\
& \leq \dfrac{w(Q_{j})}{|Q_{j}|} + \displaystyle\sum_{m=0}^{\infty}\frac{1}{2^{m-2}}[w]_{A_{1}}\operatorname{essinf}_{x \in Q_{j}}w(x) \\
& \leq \dfrac{w(Q_{j})}{|Q_{j}|} + \displaystyle\sum_{m=0}^{\infty}\frac{1}{2^{m-2}}[w]_{A_{1}}\frac{w(Q_{j})}{|Q_{j}|} \\
& = (1 + 8[w]_{A_{1}})\dfrac{w(Q_{j})}{|Q_{j}|} \leq 9[w]_{A_{1}}\dfrac{w(Q_{j})}{|Q_{j}|}.
\end{align*}
因此由 \eqref{eq:1.11.22},\eqref{eq:1.11.19},
\begin{equation*}
w(E_{2}) \leq \frac{16}{\pi}\sum_{j}9[w]_{A_{1}}w(Q_{j}) \leq C_{5}w(\Omega) \leq \frac{C_{6}}{\lambda}\| f \|_{L^{1}(w\dd x)}.
\end{equation*}
因此与 \eqref{eq:1.11.21} 合并,
\begin{equation}\tag{1.11.23}\begin{aligned}
w(\{x \in \mathbb{R};\ |Hf(x)| > \lambda,\ x \in (\Omega^{*})^{c}\}) & \leq w(E_{1}) + w(E_{2}) \\
& \leq \dfrac{C_{4} + C_{6}}{\lambda}\| f \|_{L^{1}(w\dd x)}.
\end{aligned}\label{eq:1.11.23}\end{equation}
由 \eqref{eq:1.11.20},\eqref{eq:1.11.23},
\begin{align*}
w(\{x \in \mathbb{R};\ |Hf(x)| > \lambda\}) & \leq w(\Omega^{*}) + w(\{x \in \mathbb{R};\ |Hf(x)| > \lambda,\ x \in (\Omega^{*})^{c}\}) \\
& \leq \dfrac{C_{2} + C_{4} + C_{6}}{\lambda}\| f \|_{L^{1}(w\dd x)}.
\end{align*}
由引理~\ref{lem:1.8.4},对任意 $f \in L^{1}(w\dd x)$,$Hf$ 被唯一地定义,且此估计式成立.
\end{proof}
